\documentclass[10pt,letterpaper]{article}
\usepackage{cornell-notes}

\title{Chapter 7: Next Permutation}
\author{}
\date{}

\setDocTitle{Chapter 7: Next Permutation}
\setDocSubtitle{Combinatorial Algorithms}
\setDocVersion{v1.0}
\setDocStatus{Study Companion}
\setDocOwner{Combinatorial Algorithms Cornell Notes}
\setDocAuthor{}
\setDocDate{\today}
\setDocSummary{Cornell-format study and review notes for Chapter 7: Next Permutation.}

\setCornellCollection{Combinatorial Algorithms Cornell Notes}
\setCornellUnitType{Chapter}
\setCornellUnitNumber{7}
\setCornellUnitTitle{Next Permutation}
\setCornellCompanionLabel{Study and review companion}

\hypersetup{pdftitle={Chapter 7: Next Permutation},pdfsubject={Cornell notes},pdfauthor={}}

\begin{document}
\maketitle
\begin{tcolorbox}[colback=Paper,colframe=Ink]{\small\bfseries\color{Teal}CORNELL NOTES}\par{\LARGE\bfseries Chapter 7: Next Permutation of $n$ Letters (NEXPER)}\par\medskip\textbf{Topic:} Transposition-adjacent permutation sequencing\hfill\textbf{Date:} \today\par\textbf{Study purpose:} Understand a successor obtainable from any permutation plus its sign.\end{tcolorbox}
\begin{tcolorbox}[colback=Teal!5,colframe=Teal,title=Essential question]How can all $n!$ permutations be sequenced so consecutive outputs differ by one transposition and a successor can be found without knowing its rank?\end{tcolorbox}
\begin{tcolorbox}[colback=white,colframe=Mist,title=Learning objectives]\begin{itemize}\item Interpret permutation generation as a Hamilton path problem.\item Explain the recursive list construction.\item Encode a permutation by inversion-table digits.\item Use the sign to select the fast or searching transition.\item State the bounded-average-work result and termination pattern.\end{itemize}\textbf{Prerequisites:} permutations, transpositions, inversion parity, Hamilton paths.\end{tcolorbox}
\section*{Cornell notes}
\begin{longtable}{@{}Q!{\color{Mist}\vrule width 1pt}A@{}}\cornellhead
\cue{What adjacency is required?}{The vertices are the $n!$ permutations, with an edge when one transposition converts one permutation into the other. A valid ordering is a Hamilton path in this graph. Minimum-change adjacency lets a consumer update information affected by only the exchanged positions.}
\cue{What recursive construction is used?}{Build the list for $n+1$ from transformed and alternately reversed copies of the list for $n$, appending the omitted letter. The alternating orientation ensures that the boundary between consecutive blocks is also one transposition. Induction gives uniqueness and completeness.}
\cue{What is the inversion encoding?}{For each $i$, $d_i$ counts earlier letters larger than the next letter; $0\le d_i\le i$. Prefix sums of these digits determine which digit can move without leaving its legal range. Changing one admissible digit corresponds to exchanging the affected letter with the nearest suitable value among its prefix.}
\cue{What happens for an even permutation?}{The fastest case is immediate: exchange the first two letters and mark the result odd. No encoding scan is needed. This single swap explains why half of all transitions have constant, minimal work.}
\cue{What happens for an odd permutation?}{Scan inversion digits and their cumulative parity until finding a digit that can move in the required direction: increase when the cumulative value is odd and below its upper bound, or decrease when it is even and positive. In the actual permutation, exchange the corresponding letter with the nearest smaller or larger value in its prefix. Mark the result even.}
\cue{Why can the routine start anywhere?}{The successor decision uses only the current permutation and its sign, not its position in the enumeration. Therefore any valid permutation may be supplied with the correct parity. Initialization at the identity is merely the path's standard starting point.}
\cue{Cost and termination}{The source sums the decreasing frequencies of deeper scans and obtains a bounded average amount of work per permutation, independent of $n$. The last permutation has a parity-dependent explicit pattern; the routine compares against it to lower the more-to-come flag.}
\cue{Invariants and cautions}{The array must remain a bijection of $\{1,\ldots,n\}$, and every transition must flip permutation sign because a transposition is odd. Supplying the wrong sign selects the wrong branch. Do not confuse this ordering with lexicographic order or with an adjacent-position-only method.}
\cue{Historical or legacy context}{The caller supplies and receives a logical parity flag and a completion flag. A modern iterator can validate the input permutation once, calculate its sign when starting midstream, and encapsulate the two flags.}
\cue{Retrieval practice}{\begin{enumerate}\item What graph is traversed?\item Why does every transition flip sign?\item What is $d_i$?\item Which transition is used for an even state?\item Why is rank unnecessary?\end{enumerate}}
\end{longtable}
\begin{tcolorbox}[colback=Ink!4,colframe=Ink,title=Cornell summary]
NEXPER enumerates permutations along a path whose edges are transpositions. A recursive block construction establishes that every permutation appears once and adjacent blocks meet by a single swap. The nonrecursive successor uses inversion-table information and the permutation's sign. Even states swap the first two letters immediately; odd states locate the first inversion digit that can move in the parity-prescribed direction and realize that digit change by one value-based transposition. The routine can therefore continue from any permutation when its sign is known. Half the states take the constant fast path, and the source's frequency calculation bounds average work independently of $n$. Correctness depends on preserving bijectivity, flipping sign at every call, and recognizing the explicit final pattern.
\end{tcolorbox}
\end{document}
